% ID: SEL-D0012-BC-E06
% ESTADO: cotejada. Segunda pasada visual completa y borrador compilado revisados.
% Fuente: Selectividad/matematicas_aplicadas/EXTRA_RESERVA_B_Modelo-6-Enunciados (Mat-Ap-2021).pdf
% SHA-256: 1752ec2f46b4ef2cc5d1c3346d69b79c2cc8a7aa12f6a0cd8f0d994b5bcda157
% Página: 2; bloque C; ejercicio original 6.

La probabilidad de que una persona sana se contagie de otra enferma por un
virus es del 80\% si coinciden en una reunión.
\begin{enumerate}[label=\alph*)]
  \item \textbf{(1 punto)} Si una persona enferma se reúne con dos personas
  sanas, teniendo en cuenta que contagiar a distintas personas son sucesos
  independientes entre sí, \textquestiondown{}cuál es la probabilidad de que
  se contagien las dos personas a la vez? \textquestiondown{}Cuál es la
  probabilidad de que se contagie alguna de ellas?
  \item \textbf{(1.5 puntos)} Una prueba para detectar la enfermedad da el
  resultado correcto en el 90\% de los casos cuando se le aplica a personas
  contagiadas y da falsos positivos en el 5\% de los casos cuando se aplica
  a personas sanas. Si una persona sana se reúne con una enferma y resulta
  positivo en una prueba posterior, \textquestiondown{}qué probabilidad hay
  de que se haya contagiado en la reunión?
\end{enumerate}
